\documentclass[11pt]{article}
\usepackage[margin=1in]{geometry}
\usepackage{booktabs}
\usepackage{array}
\usepackage{amsmath}
\usepackage{graphicx}
\usepackage[hidelinks]{hyperref}
\usepackage{xcolor}
\usepackage{multirow}

\title{\textbf{nous on MemoryAgentBench:\\A Paper-Faithful Evaluation of a Cognitive Memory Agent}}
\author{Timur Fatykhov\\\small\textit{with Claude (Anthropic) as evaluation engineer}}
\date{July 4, 2026}

\newcommand{\nous}{\textsf{nous}}
\newcommand{\mab}{MemoryAgentBench}

\begin{document}
\maketitle

\begin{abstract}
\noindent We evaluate \nous{}, a cognitive memory agent with episodic ingestion,
LLM-based fact extraction, and sleep-cycle consolidation, on \mab{}
(Hu et al., ICLR 2026), a benchmark that measures four memory competencies of
conversational agents: Accurate Retrieval (AR), Test-Time Learning (TTL),
Long-Range Understanding (LRU), and Conflict Resolution (CR).
We reproduce the benchmark's evaluation protocol exactly---its per-dataset answer
prompts, its per-dataset graders (including two GPT-4o judge protocols), vendored
from the benchmark's open-source repository and verified line-by-line by an
independent review---and drive a real \nous{} HTTP server through its full
production cognitive cycle for every test instance.
On the evaluated subsets, \nous{} scores
\textbf{AR 0.897} (n=232, 95\% CI [0.86, 0.94]),
\textbf{CR 0.766} (n=64),
\textbf{LRU 0.824} (n=68, CI [0.73, 0.91]),
and \textbf{TTL 0.555} (n=200, CI [0.49, 0.62]).
The first three are robustly above every method reported by the benchmark
(best published per-competency averages: 71.8\%, 29.5\%, 62.2\% respectively);
TTL sits at the top of the published range (53.9\%) but does not clearly exceed it.
We describe the testing technique, a statistical-validation pass that materially
corrected one competency's initial estimate, and the threats to validity that
qualify a direct leaderboard comparison.
\end{abstract}

\section{Introduction}

Long-running LLM agents need memory that outlives a single context window:
they must store what they were told, learn new skills from it, reason over
inputs far larger than the model context, and revise beliefs when information
changes. \mab{}~\cite{mab} operationalizes these requirements as four
competencies and evaluates them by streaming large corpora into an agent
incrementally and then asking questions against what the agent retained.

\nous{} is a cognitive AI agent whose memory subsystem combines an episodic
store (conversation episodes and summaries), a semantic store (LLM-extracted
facts with typed graph edges), full-fidelity content chunks with hybrid
(vector + lexical) retrieval, and an offline \emph{sleep} consolidation cycle
that links, deduplicates, and densifies memory between sessions. This report
documents a from-scratch evaluation harness that measures \nous{} on \mab{}
under its \emph{production} configuration, and places the results against the
methods evaluated by the benchmark's authors.

Two principles governed the effort:
\begin{enumerate}
  \item \textbf{Paper fidelity.} The score must be computed exactly as the
        benchmark computes it---same answer prompts, same graders, same judge
        models---so numbers are comparable. All grading code was vendored from
        the benchmark's MIT-licensed repository and cross-reviewed against it
        by an independent reviewer agent with no stake in the outcome.
  \item \textbf{Statistical honesty.} No above/below-field claim is made unless
        the 95\% confidence interval of the measured accuracy clears the
        best published field value. This rule changed one headline conclusion
        (Section~\ref{sec:validation}).
\end{enumerate}

\section{System Under Test}

\subsection{The \nous{} cognitive cycle}
Each benchmark instance exercises the complete production write path:

\begin{enumerate}
  \item \textbf{Ingestion.} The instance's context is split into chunks of at
        most 32{,}000 characters (conversational corpora are packed
        turn-by-turn into chunks of the same size, preserving every turn) and
        streamed to the agent as consecutive chat turns in one session.
  \item \textbf{Session close.} The session is closed, triggering the
        durable-write chain: an episode summarizer, an LLM fact extractor
        that emits typed facts, and full-fidelity chunk embedding.
  \item \textbf{Sleep consolidation.} One or more sleep cycles run offline
        consolidation: fact deduplication, cross-episode linking,
        co-occurrence edge formation, and graph densification.
  \item \textbf{Answering.} Every question is asked in a \emph{fresh} session.
        The agent answers from its memory (retrieval-augmented context
        injection plus agentic recall tools); the question is framed with the
        benchmark's own per-dataset prompt template.
\end{enumerate}

\subsection{Configuration}
\nous{} ran under its live production configuration (Claude Sonnet~5 as the
reasoning model; hybrid chunk search; typed-edge extraction; cross-encoder
re-ranking; date-aware retrieval; adaptive thinking) with exactly three
eval-hygiene overrides, all disabling background self-improvement features
irrelevant to memory measurement: the heartbeat timer, the scheduler, and
rubric outcome detection. A 117-knob snapshot of the configuration is pinned
in the harness repository for reproducibility.

Two run profiles were used. The \emph{full-fidelity} profile (unbounded
episode summarization, three sleep cycles) is used for moderate contexts.
For multi-megabyte contexts a \emph{big-context} profile bounds the episode
summarizer to 8 chunks, runs one sleep cycle, and paces ingestion turns
(15\,s); this profile eliminated upstream rate-limit stalls and ingest
timeouts entirely (0 errors across all final runs, including 2.3M-character
documents).

\subsection{Isolation}
Every instance gets a fresh \nous{} server process, a unique agent identity,
and a dedicated evaluation database; nothing carries over between instances.
Control probes confirmed the isolation: on Conflict Resolution, the identical
questions asked of an agent \emph{before} ingestion score 0.000, so the
measured accuracy is attributable to memory, not to model priors.

\section{Benchmark and Datasets}

\mab{} streams each dataset's context into the agent and then asks held-out
questions. Table~\ref{tab:datasets} lists the evaluated subsets.

\begin{table}[h]
\centering\small
\begin{tabular}{llrrl}
\toprule
Competency & Source & Context size & Questions used & Grader \\
\midrule
\multirow{4}{*}{CR}
 & factconsolidation\_sh (6k--262k)  & 26k--1.1M chars & $4\times8$ & substring (SQuAD-norm.) \\
 & factconsolidation\_mh (6k--262k)  & 26k--1.1M chars & $4\times8$ & substring (SQuAD-norm.) \\
 & \multicolumn{4}{l}{\footnotesize single-hop (sh) = updated fact; multi-hop (mh) = compose two updated facts} \\[2pt]
\multirow{6}{*}{AR}
 & eventqa\_65536                     & $\sim$285k chars & 200 (5 ctx $\times$ 40) & all-elements recall \\
 & eventqa\_131072                    & $\sim$570k chars & 8  & all-elements recall \\
 & eventqa\_full                      & $\sim$2.3M chars & 8  & all-elements recall \\
 & ruler\_qa1\_197K / qa2\_421K       & 0.9M / 1.9M chars & 8+8 & substring \\
 & longmemeval\_s*                    & 1.6M chars (1282 turns) & 8 & GPT-4o yes/no judge \\[2pt]
\multirow{2}{*}{TTL}
 & icl\_banking77 / clinic150 / nlu   & $\sim$400k chars each & $3\times40$ & parsed-label exact match \\
 & icl\_trec\_coarse / trec\_fine     & $\sim$400k chars each & $2\times40$ & parsed-label exact match \\[2pt]
\multirow{2}{*}{LRU}
 & detective\_qa (10 novels)          & $\sim$400k chars each & 68 & substring \\
 & infbench\_sum\_eng                 & 1.75M chars & 1 & GPT-4o 3-call f1 judge \\
\bottomrule
\end{tabular}
\caption{Evaluated datasets. ``Context size'' is the raw document streamed into
memory. The TTL \emph{recsys} sub-dataset could not be graded faithfully: its
grader requires an entity-id mapping file that is not shipped in the
benchmark's public repository, so TTL here is the in-context-learning family
only.}
\label{tab:datasets}
\end{table}

\section{Testing Technique}

\subsection{Paper-faithful grading}
Each sub-dataset is graded exactly as in the benchmark's \texttt{post\_process}
dispatch:

\begin{itemize}
  \item \textbf{Substring match} (ruler, detective, factconsolidation): the
        SQuAD-style normalization (lowercase, strip punctuation and articles,
        collapse whitespace), then gold-in-prediction containment, max over
        golds.
  \item \textbf{EventQA recall}: \emph{raw} case-insensitive containment of
        every gold element (\texttt{el.lower() in pred.lower()}); the answer is
        correct only if all elements are present. An initial implementation
        that reused the SQuAD normalization here was flagged by independent
        review as strictly more lenient than the paper (punctuation and
        article stripping create false positives) and was corrected before any
        reported run.
  \item \textbf{ICL exact match}: parse the label the prompt requests
        (\texttt{label:~X}), then exact match on the bare label. This is a
        deliberate, conservative interpretation: the benchmark's own ICL
        scoring parses with a different prefix and reports multiple metrics of
        which the substring variant false-positives on digit collisions
        (gold \texttt{5} matches \texttt{label: 15}); our parser cannot.
  \item \textbf{LongMemEval judge}: the benchmark's GPT-4o yes/no judge with
        its verbatim per-question-type prompts (including the
        temporal-reasoning off-by-one tolerance and the abstention protocol
        for \texttt{\_abs} questions).
  \item \textbf{Summarization judge} (InfBench-Sum): the benchmark's three-call
        GPT-4o protocol scoring fluency $f\in\{0,1\}$, key-point recall $r$,
        and sentence-level precision $p$, combined as
        $\mathrm{f1} = f \cdot \frac{2rp}{r+p}$.
\end{itemize}

Answer prompts are likewise the benchmark's own per-dataset templates (e.g.,
the CR template that states the serial-number update rule with few-shot
examples; the eventQA template that frames ``what happened next''). On
identical memory, switching from a generic answer prompt to the benchmark's CR
template moved CR accuracy by +28 percentage points---prompt fidelity is not
optional if numbers are to be compared.

\subsection{Independent review}
All vendored grading and prompting code was audited by a reviewer agent with
fresh context against the benchmark's repository, file by file. The review
confirmed byte-fidelity of the normalizers, judge prompts, dispatch tables,
and the summarization formula; it found exactly one fidelity bug (the eventQA
normalization above), which was fixed and regression-tested before the
reported runs. The harness carries 125 unit tests covering graders, judges,
chunking, persistence, and isolation.

\subsection{Statistical validation}
\label{sec:validation}
Initial runs used 8 questions per source. A dedicated validation pass then
grew the sample until each competency's 95\% CI either cleared the best
published field value or demonstrably did not:

\begin{itemize}
  \item \textbf{TTL corrected downward.} At $n{=}8$/source TTL measured 0.650;
        at $n{=}40$/source ($n{=}200$) it fell to \textbf{0.555}
        (CI [0.486, 0.624]), with one source falling from 0.625 to 0.450.
        The small sample had been an optimistic draw. TTL is therefore
        reported as \emph{at} the field ceiling, not above it.
  \item \textbf{LRU confirmed.} Detective-QA at all 10 novel contexts
        ($n{=}68$) held at \textbf{0.824} versus 0.833 at $n{=}6$; its lower
        CI bound (0.733) clears the best field value (0.622).
  \item \textbf{AR firmed upward.} EventQA-65536 across 5 contexts at
        $n{=}200$ rose to 0.910 versus 0.875 at $n{=}24$, consolidating AR at
        \textbf{0.897} (CI [0.857, 0.936]).
\end{itemize}

The general lesson---small-sample memory-eval estimates run optimistic; grow
$n$ until the CI clears the comparison band---is recorded as a standing
methodological rule for this project. All results persist per-question to
append-only JSONL as each instance completes, so partial runs are never lost
and every reported number is recomputed from the persisted records.

\section{Results}

\subsection{\nous{} scores}

\begin{table}[h]
\centering
\begin{tabular}{lrrll}
\toprule
Competency & Score & $n$ & 95\% CI & Best published avg.\ / verdict \\
\midrule
Accurate Retrieval        & \textbf{0.897} & 232 & [0.857, 0.936] & 0.718 --- \textbf{above, robust} \\
Conflict Resolution       & \textbf{0.766} & 64  & [0.66, 0.87]   & 0.295 --- \textbf{above, robust} \\
Long-Range Understanding  & \textbf{0.824} & 68  & [0.733, 0.914] & 0.622 --- \textbf{above, robust} \\
Test-Time Learning        & \textbf{0.555} & 200 & [0.486, 0.624] & 0.539 --- at ceiling \\
\bottomrule
\end{tabular}
\caption{\nous{} on \mab{} (paper-faithful grading, production configuration,
0 execution errors on all final runs). LRU is detective-QA; the auxiliary
InfBench-Sum summarization judge scored f1 = 0.592 on a single 1.75M-char book
(unfirmed, $n{=}1$). ``Best published avg.'' is the strongest per-competency
average across all methods in the benchmark's results table.}
\label{tab:results}
\end{table}

Per-source detail: CR single-hop 0.906 vs.\ multi-hop 0.625 (the benchmark
reports at most $\sim$7\% for multi-hop across all its methods); AR eventqa
0.910/0.750/0.875 (65k/131k/full), ruler 0.875/0.750 (197K/421K), longmemeval
0.875; TTL banking77 0.525, clinic150 0.500, nlu 0.425, trec-coarse 0.875,
trec-fine 0.450; LRU detective 56/68.

\subsection{Comparison with the field}

Table~\ref{tab:field} reproduces the per-competency averages of representative
methods from the benchmark's results table (its Table~3; the fourth competency
is named Selective Forgetting there and Conflict Resolution in earlier
versions---same datasets).

\begin{table}[h]
\centering
\begin{tabular}{llrrrr}
\toprule
Method & Class & AR & TTL & LRU & CR/SF \\
\midrule
GPT-4o-mini (full context) & long-context & 49.2 & 48.6 & 46.2 & 25.0 \\
GPT-4.1-mini (full context) & long-context & 71.8 & 46.2 & 49.1 & 20.5 \\
Claude-3.7-Sonnet (full context) & long-context & 59.7 & \textbf{53.9} & \textbf{62.2} & 22.5 \\
BM25 & RAG & 60.5 & 44.5 & 35.6 & 25.5 \\
Text-Embed-3-Large & RAG & 54.6 & 44.3 & 37.3 & 16.0 \\
HippoRAG-v2 & RAG & 65.1 & 35.8 & 36.2 & \textbf{29.5} \\
Mem0 & agentic memory & 32.6 & 21.2 & 20.7 & 10.0 \\
GraphRAG & structure-aug. & 40.9 & 24.8 & 19.9 & 8.0 \\
MIRIX & agentic memory & 47.5 & 24.1 & 25.4 & 8.0 \\
\midrule
\textbf{\nous{}} (this work) & cognitive memory & \textbf{89.7} & 55.5 & \textbf{82.4} & \textbf{76.6} \\
\bottomrule
\end{tabular}
\caption{Per-competency averages (\%). Field rows are from the benchmark's
published results; bold marks the best value per column among field methods,
and \nous{} where it exceeds all of them. Comparison caveats in
Section~\ref{sec:threats} apply---in particular the base-model mismatch and
the subset sizes.}
\label{tab:field}
\end{table}

Three observations:
\begin{enumerate}
  \item \textbf{Conflict Resolution is the standout.} \nous{} scores 0.766
        where no published method exceeds 0.295, and 0.625 on multi-hop
        conflict composition where the field is at most $\sim$0.07. We
        attribute this to \nous{}'s explicit fact store with update semantics:
        conflicting facts are extracted, linked, and consolidated during
        sleep, rather than left as contradictory passages for retrieval to
        stumble over.
  \item \textbf{Retrieval and long-range understanding benefit from
        full-fidelity chunk memory.} AR 0.897 and LRU 0.824 exceed both the
        RAG family and the long-context family, including on 1.6--2.3M-char
        inputs that no context window holds.
  \item \textbf{Test-time learning is the honest gap.} TTL 0.555 edges the
        best field value (0.539) but the CI straddles it. Learning a
        classification skill from thousands of in-context examples is not
        what an episodic/semantic memory is optimized for; accuracy here
        tracks how many labeled exemplars retrieval can surface per query.
\end{enumerate}

\section{Threats to Validity}
\label{sec:threats}

\begin{enumerate}
  \item \textbf{Base-model mismatch.} \nous{} reasons with Claude Sonnet~5;
        the benchmark's agentic-memory rows mostly use GPT-4o-mini, and its
        strongest long-context rows use GPT-4.1-mini and Claude-3.7-Sonnet.
        Part of \nous{}'s margin is attributable to the stronger base model.
        The long-context Claude-3.7 row (best field TTL and LRU) is the
        fairest anchor, and \nous{} still exceeds it on AR (+30pp), LRU
        (+20pp), and CR (+54pp).
  \item \textbf{Subsets, not the full benchmark.} Question counts per source
        range from 8 to 200 (Table~\ref{tab:datasets}); the giants
        (eventqa-full, ruler-421K, longmemeval) were measured at one context
        each. CIs in Table~\ref{tab:results} reflect this. recsys (TTL) is
        excluded for lack of the grader's data file; InfBench-Sum is a single
        book.
  \item \textbf{Configuration provenance.} The production configuration was
        improved using failures observed on this benchmark's CR subset in an
        earlier phase (write-path fidelity, hybrid chunk retrieval, recall
        depth). The evaluation is therefore best read as ``\nous{} after its
        memory bugs were fixed,'' not as a blind first attempt. No knob was
        tuned per-dataset, and the same configuration serves all four
        competencies.
  \item \textbf{Ingestion transport differs from the paper's protocol.} The
        benchmark feeds context in chunks of 512 tokens for the doc-QA,
        LongMemEval, and conflict-resolution tasks and 4096 tokens elsewhere
        (uniformly 4096 for the commercial memory agents, for cost); our
        harness streams 32{,}000-character turns, and \nous{} re-chunks
        internally for retrieval, so transport granularity is not the
        retrieval granularity. The baseline is therefore paper-faithful in
        prompts, graders, judges, and question sets, but \emph{not} an exact
        reproduction of the paper's ingestion transport. No answer-bearing
        text was lost to this difference: ingest capacity was
        80 $\times$ 32k chars (2.56M), and a post-hoc audit of all 34
        evaluated instances at the exact run settings confirmed
        \textbf{zero dropped chunks} (largest instance: 2.31M chars, 73/80
        chunks; the 1{,}282-turn LongMemEval conversation packed to 50 chunks
        with every turn retained). Runs now stamp
        \texttt{chunks\_sent}/\texttt{chunks\_truncated} on every persisted
        result row, and truncated instances are excluded from headline
        aggregates by construction.
  \item \textbf{Benchmark prompts embed task rules.} The benchmark's own CR
        and eventQA templates tell the agent the update rule / task framing.
        We use them because the field numbers use them; they help any
        system.
  \item \textbf{Judge sharing.} LongMemEval and InfBench-Sum scores depend on
        GPT-4o judges, as in the benchmark.
\end{enumerate}

\section{Conclusion}

Under the benchmark's exact evaluation protocol and its production
configuration, \nous{} is robustly above every method published with
\mab{} on three of four memory competencies---Accurate Retrieval (0.897),
Conflict Resolution (0.766, with a 2.6$\times$ margin over the best published
average), and Long-Range Understanding (0.824)---and statistically tied with
the best published value on Test-Time Learning (0.555). The conflict-resolution
result in particular suggests that an explicit extract-link-consolidate memory
architecture addresses a failure mode that both long-context and RAG systems
share. The evaluation harness, vendored graders, pinned configuration, and
per-question result records are available in the project repository for
reproduction.

\begin{thebibliography}{9}
\bibitem{mab}
Y.~Hu et al.,
``Evaluating Memory in LLM Agents via Incremental Multi-Turn Interactions''
(\mab{}), ICLR 2026. \texttt{github.com/HUST-AI-HYZ/MemoryAgentBench},
arXiv:2507.05257.
\end{thebibliography}

\end{document}
